% =============================================================================
% layout.tex —— 版式层：页标题、边距、页脚与页码、过渡页、目录与列表模板
% -----------------------------------------------------------------------------
% 依赖：colors.tex（字体颜色）、style.tex（\ink@margin）。
% 本文件是「改版面观感」的主要入口。
% =============================================================================

% -----------------------------------------------------------------------------
% 1. 页标题与边距
% -----------------------------------------------------------------------------
% 标题、横线和正文共享同一组左右边距，只需在 style.tex 里改 \ink@margin。
\setbeamersize{text margin left=\ink@margin,text margin right=\ink@margin}

% 标题下方蓝色横线的粗细
\newcommand{\ink@toprule}{0.7pt}

% 自定义页标题模板：标题（+可选副标题）→ 蓝色细横线 → 下方留白。
% 用 \hrule 而不是 \hline，避免额外的表格上下文；
% 顶部 0.35cm 与底部 0.25cm 的留白决定正文起始位置。
\setbeamertemplate{frametitle}{%
  \nointerlineskip
  \begin{beamercolorbox}[wd=\textwidth,sep=0pt]{frametitle}
    \vskip0.35cm
    \usebeamerfont{frametitle}\strut\insertframetitle\strut\par
    \ifx\insertframesubtitle\empty\else
      \vskip0.2em
      {\usebeamerfont{framesubtitle}\usebeamercolor[fg]{framesubtitle}\insertframesubtitle\par}%
    \fi
    \vskip0.12cm
    {\color{AcademicBlue}\hrule height\ink@toprule}
    \vskip0.25cm
  \end{beamercolorbox}}

% 多页 frame 自动续页时，在标题后追加序号，例如「（2）」
\setbeamertemplate{frametitle continuation}{%
  \ifnum\insertcontinuationcount>1\relax\space(\insertcontinuationcount)\fi}

% 标题、章节页用衬线粗体（中文即思源宋体 Bold）
\setbeamerfont{title}{family=\rmfamily,series=\bfseries}
\setbeamerfont{frametitle}{family=\rmfamily,series=\bfseries}
\setbeamerfont{framesubtitle}{family=\sffamily,series=\mdseries,shape=\upshape}
\setbeamerfont{section title}{family=\rmfamily,series=\bfseries}

% 封面保留原有字段与字号，空的作者/单位不占位，使标题与日期更紧凑。
\setbeamertemplate{title page}{%
  \begingroup
  \centering\setlength{\parskip}{0pt}%
  % 原生圆角色盒，无阴影；宽度随页面缩放，长标题可自动换行。
  \noindent\makebox[\textwidth][c]{%
    \begin{minipage}{0.84\textwidth}
      \begin{beamercolorbox}[wd=\linewidth,sep=0.9em,center,rounded=true,shadow=false]{title panel}
        \usebeamerfont{title}\inserttitle\par
      \end{beamercolorbox}
    \end{minipage}}\par
  \ifx\insertsubtitle\empty\else
    \vspace{0.5em}%
    {\usebeamerfont{subtitle}\usebeamercolor[fg]{subtitle}\insertsubtitle\par}%
  \fi
  \ifx\insertauthor\empty\else
    \vspace{1em}{\usebeamerfont{author}\usebeamercolor[fg]{author}\insertauthor\par}%
  \fi
  \ifx\insertinstitute\empty\else
    \vspace{0.5em}{\usebeamerfont{institute}\usebeamercolor[fg]{institute}\insertinstitute\par}%
  \fi
  \vspace{1.4em}%
  {\usebeamerfont{date}\usebeamercolor[fg]{date}\insertdate\par}%
  {\usebeamercolor[fg]{titlegraphic}\inserttitlegraphic\par}%
  \endgroup}

% 关闭右下角导航图标；文献表不显示条目图标
\setbeamertemplate{navigation symbols}{}
\setbeamertemplate{bibliography item}{}
% 图表编号
\setbeamertemplate{caption}[numbered]
% 图表采用相同的题注间距；图表专属来源在对应 figure/table 内排版。
\setlength{\abovecaptionskip}{0.45em}
\setlength{\belowcaptionskip}{0.35em}

% -----------------------------------------------------------------------------
% 2. 列表模板
% -----------------------------------------------------------------------------
\setbeamertemplate{enumerate item}{\insertenumlabel.}
\setbeamertemplate{enumerate subitem}{\insertenumlabel.\insertsubenumlabel.}
\setbeamertemplate{itemize item}{\small\textbullet}
\setbeamertemplate{itemize subitem}{\textendash}

% 单栏短要点：适度下移起点并拉开条目，不影响目录、双栏与块内列表。
\newenvironment{inkpoints}{%
  \par\vspace{0.12cm}%
  \begin{enumerate}\setlength{\itemsep}{0.55em}%
}{\end{enumerate}}

% -----------------------------------------------------------------------------
% 3. 页脚与页码（含附录独立编号）
% -----------------------------------------------------------------------------
% 正文页码：「当前 / 总」。沿用 Beamer 内置的 totalframenumber 模板，
% 其内容等价于原来的 \insertframenumber\,/\,\inserttotalframenumber。
%
% 附录页码：独立编号「A-序号 / A-总数」。
%   \insertframenumberinappendix   附录内第几页（从 1 开始）
%   \insertappendixframenumber     附录总页数
% 这两个计数器由 Beamer 在 \appendix 时维护，并借助 .nav 辅助文件跨轮传递，
% 页码由 Beamer 维护；在进入附录时切换模板，不在页脚中读取内部模式条件。
\defbeamertemplate*{page number in head/foot}{inkappendixframenumber}{%
  A-\insertframenumberinappendix\,/\,A-\insertappendixframenumber}

% 当前使用的页码模板名；\appendix 会把它切换成附录模板。
% 使用 \gdef 是因为 \appendix 在正文中执行，需要全局生效。
\def\ink@pagenumberstyle{totalframenumber}

% 统一的页码输出：右对齐 + 当前模板；字号和颜色由内部样式模块设定。
% \setbeamertemplate 在这里位于每个页脚的分组内部，作用范围仅限本次排版；
% 由于模板名是全局的，切换一次即对所有后续页面稳定生效。
\newcommand{\ink@pagenumber}{%
  \setbeamertemplate{page number in head/foot}[\ink@pagenumberstyle]%
  \usebeamerfont{page number in head/foot}%
  \usebeamercolor[fg]{page number in head/foot}\hfill
  \usebeamertemplate***{page number in head/foot}\par}

% 包装 \appendix：先执行 Beamer 原本的附录处理（进入附录模式、写入计数器），
% 再把页码模板切换到附录版本。这样附录页脚显示 A-当前 / A-总数，
% 而正文始终保持「当前 / 总」不变。
\let\ink@appendixoriginal\appendix
\renewcommand{\appendix}{%
  \ink@appendixoriginal
  \global\def\ink@pagenumberstyle{inkappendixframenumber}}

% -----------------------------------------------------------------------------
% 4. 章节与小节过渡页
% -----------------------------------------------------------------------------
% 写在 frame 之外的 \section / \subsection 会自动插入一张居中过渡页，
% 无需手写额外 frame。两类过渡页都不显示页码（noframenumbering），
% 因此不会占用正文页码，也不受附录编号影响。
% 带星号的 \section* / \subsection* 不触发过渡页。

% 章节页：水平垂直居中，显示「编号 + 标题」，例如「1. 研究背景」
\AtBeginSection[]{%
  \begin{frame}[plain,c,noframenumbering]
    \centering
    {\usebeamerfont{section title}\usebeamercolor[fg]{section title}%
      {\color{AcademicBlue}\insertsectionnumber.}\enspace\insertsection\par}
  \end{frame}
}

% 小节页：上下结构。上方灰色「章节号 + 章节名」，下方较大的「章节号.小节号 + 小节名」
\AtBeginSubsection[]{%
  \begin{frame}[plain,c,noframenumbering]
    \centering
    {\normalsize\color{AcademicGray}%
      \insertsectionnumber.\enspace\insertsection\par}
    \vspace{0.8em}
    {\usebeamerfont{section title}\usebeamercolor[fg]{section title}%
      {\color{AcademicBlue}\insertsectionnumber.\insertsubsectionnumber}\enspace\insertsubsection\par}
  \end{frame}
}
